\documentclass{axiobac}

\annee{1990}
\mois{Juin}
\session{contrôle}
\epreuve{Mathématiques}
\sectionbac{Math-Sciences \& Math-Technique}
\sections{mathematiques, sciences-techniques}
\duree{4 heures}
\coefficient{5}
\ministere{Ministère de l'Éducation, de l'Enseignement Supérieur et de la Recherche Scientifique}

\begin{document}
\entete

%----------------------------------------------------------------------------------------
%   EXERCICE 1
%----------------------------------------------------------------------------------------
\begin{exercice}{4}
\titre{Un dé truqué lancé trois fois}
\chapitres{Probabilités}

On lance un dé cubique dont les faces sont numérotées de $1$ à $6$. Ce dé est truqué de façon qu'à chaque jet la probabilité d'obtenir un nombre pair soit égale au double de celle d'obtenir un nombre impair.

\begin{enumerate}
    \item On lance le dé une fois et on désigne par $p$ la probabilité d'obtenir un nombre impair.
    \begin{enumerate}
        \item Calculer $p$.
        \item En déduire la probabilité d'obtenir le nombre $1$ et celle d'obtenir le nombre $2$.
    \end{enumerate}
    \item On lance le dé trois fois. Calculer la probabilité de chacun des événements suivants :
    \begin{enumerate}
        \item on obtient deux fois et deux fois seulement le nombre $2$ ;
        \item on obtient au moins deux fois le nombre $2$.
    \end{enumerate}
\end{enumerate}

\end{exercice}

%----------------------------------------------------------------------------------------
%   EXERCICE 2
%----------------------------------------------------------------------------------------
\begin{exercice}{4}
\titre{Similitude, homothétie et rotation en complexes}
\chapitres{Similitudes, Nombres Complexes}

Le plan $P$ est rapporté à un repère orthonormé direct $(O, \vec{u}, \vec{v})$.

\begin{enumerate}
    \item Soit $g$ l'application de $P$ dans $P$ qui, au point $M$ d'affixe $z$, associe le point $M'$ d'affixe $z'$ tel que
    \[
    z' = \left(1 - i\sqrt{3}\right) z + (1 + i)\sqrt{3} .
    \]
    Déterminer la nature et les éléments caractéristiques de $g$.
    \item Soit, dans le plan $P$, l'homothétie $h$ de centre $\Omega(1, -1)$ et de rapport $\dfrac{1}{2}$. On pose $f = h \circ g$. \\
    Déterminer la nature et les éléments caractéristiques de l'application $f$.
    \item Soit $D$ la droite dont une équation est $y = x$. \\
    Déterminer et construire l'image de $D$ par l'application $f$.
\end{enumerate}

\end{exercice}

%----------------------------------------------------------------------------------------
%   PROBLÈME
%----------------------------------------------------------------------------------------
\begin{probleme}{12}

Soit $P$ le plan rapporté à un repère orthonormé $(O, \vec{i}, \vec{j})$ et $\Delta$ la droite d'équation : $y = x$.

\partie{A}
\titre{La fonction $\ln(1+2x)$, sa réciproque, une aire et une suite récurrente}
\chapitres{Fonction logarithme népérien, Fonctions réciproques, Intégrales, Suites réelles}

Soit $f$ la fonction numérique à variable réelle définie par $f(x) = \operatorname{Log}(1 + 2x)$ et $\mathcal{C}$ sa courbe représentative dans le repère $(O, \vec{i}, \vec{j})$.

\begin{enumerate}
    \item Étudier les variations de $f$.
    \item On désigne par $\varphi$ la fonction définie par : $\varphi(x) = f(x) - x$.
    \begin{enumerate}
        \item Montrer que l'équation $\varphi(x) = 0$ admet deux solutions $0$ et $\alpha$ avec $1 < \alpha < 2$.
        \item Déterminer les positions relatives de $\mathcal{C}$ et $\Delta$.
        \item Tracer $\mathcal{C}$ et $\Delta$ dans le repère $(O, \vec{i}, \vec{j})$.
    \end{enumerate}
    \item
    \begin{enumerate}
        \item Montrer que $f$ admet une fonction réciproque $f^{-1}$ dont on tracera la courbe $\mathcal{C}'$ dans le même repère.
        \item Calculer $f^{-1}(x)$ pour tout réel $x$.
        \item Calculer, en fonction de $\alpha$, l'aire du domaine limité par les courbes $\mathcal{C}$, $\mathcal{C}'$ et les droites d'équations respectives $x = 0$ et $x = \alpha$.
    \end{enumerate}
    \item On définit la suite $(U_n)$ par :
    \[
    \begin{cases}
        U_0 = 1 \\
        U_{n+1} = f(U_n), \quad n \in \mathbb{N}
    \end{cases}
    \]
    \begin{enumerate}
        \item Montrer que : $\forall n \in \mathbb{N}$, $1 \leqslant U_n < \alpha$.
        \item Montrer que $(U_n)$ est strictement croissante.
        \item En déduire que $(U_n)$ est convergente et trouver sa limite.
    \end{enumerate}
\end{enumerate}

\partie{B}
\titre{Les fonctions $\ln(1+nx)$ : point fixe et comportement asymptotique}
\chapitres{Fonction logarithme népérien}

Dans cette partie, $n$ désigne un entier naturel tel que $n \geqslant 2$. \\
Soit $f_n$ la fonction numérique à variable réelle définie par : $f_n(x) = \operatorname{Log}(1 + nx)$ et $\mathcal{C}_n$ sa courbe représentative dans le repère $(O, \vec{i}, \vec{j})$.

\begin{enumerate}
    \item Montrer que : $\forall x \in ]1, +\infty[$, $1 - \dfrac{1}{x} < \operatorname{Log} x < \dfrac{1}{2}\left(x - \dfrac{1}{x}\right)$.
    \item Montrer que l'équation $f_n(x) - x = 0$ admet deux solutions (on notera $\alpha_n$ la solution non nulle).
    \item
    \begin{enumerate}
        \item Montrer que : $\operatorname{Log} n < \alpha_n < 2 \operatorname{Log} n$.
        \item En déduire que : $\operatorname{Log} n < \alpha_n < \operatorname{Log}(1 + 2n \operatorname{Log} n)$.
        \item Calculer $\displaystyle\lim_{n \to +\infty} \frac{\alpha_n}{\operatorname{Log} n}$.
    \end{enumerate}
    \item Soit $\Phi$ l'application du plan $P$ dans lui-même qui, à tout point $M$, associe le point $M'$ barycentre des points $H$ et $M$ affectés respectivement des coefficients $(n - 2)$ et $2$, où $H$ désigne le projeté orthogonal de $M$ sur $(O, \vec{j})$.
    \begin{enumerate}
        \item Déterminer la nature de $\Phi$ et préciser ses éléments caractéristiques.
        \item Montrer que $\Phi(\mathcal{C}_2) = \mathcal{C}_n$.
        \item Construire $\mathcal{C}_4$ dans le même repère $(O, \vec{i}, \vec{j})$.
    \end{enumerate}
\end{enumerate}

\end{probleme}

\end{document}
